% Infraestrutura tipográfica da edição portuguesa de SGA 3, VIII--IX.
% Este arquivo define somente a composição; não é fonte do texto matemático.
\documentclass[11pt,a4paper,oneside,openany]{book}
\usepackage[T1]{fontenc}
\usepackage[utf8]{inputenc}
\usepackage{lmodern}
\usepackage{amsmath,amssymb,amsthm,mathrsfs,mathtools}
\usepackage[a4paper,margin=25mm,headheight=14pt,headsep=7mm,footskip=12mm]{geometry}
\usepackage{microtype}
\usepackage{array,longtable,enumitem}
\usepackage{tikz-cd}
\usetikzlibrary{arrows,arrows.meta,positioning}
\usepackage{fancyhdr}
\usepackage[unicode,hidelinks,hypertexnames=false]{hyperref}
\usepackage{bookmark}
\IfFileExists{brazilian.ldf}{\usepackage[brazilian]{babel}}{%
  \usepackage[provide=*]{babel}
  \babelprovide[import,main]{brazilian}}
% Com brazilian.ldf, o babel torna " um caractere ativo; a biblioteca babel do TikZ
% protege os rótulos "..." dos diagramas tikz-cd (27 nov. 2026).
\usetikzlibrary{babel}
\hypersetup{pdftitle={SGA 3, Exposições VIII e IX — texto de 1970},
  pdfauthor={A. Grothendieck},pdfsubject={Tradução em português brasileiro}}
\setlength{\parindent}{1.5em}
\setlength{\parskip}{0pt plus .3pt}
\setlength{\emergencystretch}{2em}
\setlist{topsep=.5\baselineskip,itemsep=.25\baselineskip,parsep=0pt}
\allowdisplaybreaks[1]
\pagestyle{fancy}
\fancyhf{}
\fancyhead[L]{\small\nouppercase{\leftmark}}
\fancyfoot[C]{\thepage}
\renewcommand{\headrulewidth}{0pt}
\fancypagestyle{plain}{\fancyhf{}\fancyfoot[C]{\thepage}\renewcommand{\headrulewidth}{0pt}}
\raggedbottom

% As chamadas conservam os nomes e identificadores dos arquivos recebidos.
% Remissões a exposições ausentes conservam o texto, sem destino fictício.
\makeatletter
\newcommand{\SGASetTarget}[1]{\phantomsection\label{#1}}
\newcommand{\SGARefTarget}[1]{%
  \@ifnextchar\chapter{\SGAChapterTarget{#1}}{\SGASetTarget{#1}}}
\def\SGAChapterTarget#1\chapter*#2{\chapter*{#2}\SGASetTarget{#1}}
\newcommand{\SGAInlineRefTarget}[1]{\SGASetTarget{#1}}
\newcommand{\SGARefLink}[2]{\@ifundefined{r@#1}{#2}{\hyperref[#1]{#2}}}
\makeatother

\newcommand{\SGAFrontCover}[7]{%
  \begin{titlepage}\centering
  {\Large #1\par}\vspace{1em}{\large #2\par}\vspace{1em}
  {#3\par}\vspace{2.5em}{\Huge #4\par}\vfill
  {\LARGE\bfseries #5\par}\vfill
  {\Large #6\par}\vspace{1em}{#7\par}
  \end{titlepage}}
\newcommand{\SGACopyrightPage}[1]{%
  \clearpage\thispagestyle{empty}\vspace*{\fill}
  {\small\noindent #1\par}\clearpage}
\newcommand{\SGAInnerTitle}[5]{%
  \clearpage\thispagestyle{empty}\begin{center}
  {\Large #1\par}\vspace{1em}{\large #2\par}\vspace{2em}
  {\LARGE\bfseries #3\par}\vspace{2em}{\large #4\par}\vfill
  {\Large #5\par}\end{center}\clearpage}
\newenvironment{SGAContentsPage}[1]{%
  \clearpage\thispagestyle{plain}\begingroup\small
  \setlength{\parindent}{0pt}\setlength{\parskip}{0pt}
  \if\relax\detokenize{#1}\relax\else
    \begin{center}\large\bfseries #1\end{center}\medskip\fi
}{\par\endgroup\clearpage}
\newcommand{\SGAContentsExpose}[3]{%
  \par\addvspace{.8em}\noindent\textbf{#1}\par
  \begingroup\raggedright\noindent #2\nobreak\hfill #3\par\endgroup\smallskip}
\newcommand{\SGAContentsEntry}[3]{%
  \par\begingroup\hangindent=2em\hangafter=1
  \noindent\makebox[2em][l]{#1}#2\nobreak\hfill #3\par\endgroup}
\newcommand{\SGAContentsSubentry}[3]{%
  \par\begingroup\leftskip=2em\hangindent=2em\hangafter=1
  \noindent\makebox[2em][l]{#1}#2\nobreak\hfill #3\par\endgroup}
\newcommand{\SGAContentsAppendix}[3]{%
  \par\addvspace{.8em}\noindent\textbf{#1} #2\nobreak\hfill #3\par\smallskip}
\newcommand{\SGAContentsIndex}[2]{%
  \par\addvspace{.8em}\noindent\textbf{#1}\nobreak\hfill #2\par}
